\section{Separating pairs in a common bag}
\label{sec:perturbations}

Fix an integer $k\geq 2$. Throughout this section, $(V,d)$ is the
shortest-path metric of a finite connected graph with positive edge lengths
and with a tree decomposition whose bags are nonempty and have size at most $k$. Write
$\Delta=\operatorname{diam}(V,d)$, and assume $|V|\geq 2$.
All constants in this section depend only on $k$.

We first apply the particle construction to pairs of vertices that share
a bag. At each bag, only boundedly many scales are needed to represent
its pairwise distances. Reserving a particle pair at each of these scales
will therefore fit into bounded memory. Padded partitions supply the
perturbations used at each launch; their padding guarantee will also be
used in the anchor construction of Section~\ref{sec:anchors}.

We use the following consequence of the padded-decomposition theorem for
bounded-treewidth graphs \cite[Theorem~1.2 and Definition~2.1]{FiltserEtAl}.
The diameter in its statement is
measured in the original metric, sometimes called weak diameter.

\begin{theorem}[Padded partitions]\label{thm:padded}
There are constants $\alpha\in(0,1/8]$ and $\pi\in(0,1]$ such that, for
every metric as above and every $r>0$, there is a random partition
$\mathcal P$ of $V$ satisfying
\[
  \operatorname{diam}(C)\leq r\quad(C\in\mathcal P),
  \qquad
  \mathbb P\bigl[B_d(x,\alpha r)\subseteq\mathcal P(x)\bigr]\geq\pi
  \quad(x\in V).
\]
Here $\mathcal P(x)$ denotes the part containing $x$.
\end{theorem}

To match parameters, the cited scheme has constants $\beta_t>0$ and
$\delta_t>0$ for each treewidth $t\geq2$, and at scale $r$ it gives
padding probability at least $\exp(-\beta_t\gamma)$ whenever
$0\leq\gamma\leq\delta_t$.  Put $t_{\max}=\max\{2,k-1\}$ and choose
\[
 \alpha=\min\bigl\{1/8,\min_{2\leq t\leq t_{\max}}\delta_t\bigr\},
 \qquad
 \pi=\min_{2\leq t\leq t_{\max}}\exp(-\beta_t\alpha).
\]
These are positive constants depending only on the bag-size bound $k$.

For completeness, the treewidth-one case also follows without interpreting
the source's logarithmic bound at one. Attach a positive-length triangle
at one vertex. The enlarged graph has treewidth two and preserves all
original distances. Restricting its partition parts to the original
vertices preserves both weak diameter and the padding event. For a fixed
bound $k$, take common positive constants over the finitely many possible
treewidth values.

We fix $\alpha$ and $\pi$ from Theorem~\ref{thm:padded} for the rest of
the proof. For a subset $C\subseteq V$, we interpret
$d(z,V\setminus C)=+\infty$ when $C=V$; this expression will always be
clipped by a finite constant.

\begin{lemma}[Disjoint nonnegative bumps]\label{lem:disjoint-bumps}
Let $E_1,\ldots,E_m$ be pairwise disjoint subsets of a metric space.
For each $i$, let $\psi_i$ be a nonnegative $1$-Lipschitz function
vanishing outside $E_i$, and let $u_i\in[0,1]$. Then
$\sum_i u_i\psi_i$ is $1$-Lipschitz.
\end{lemma}

\begin{proof}
Within one set $E_i$, the assertion follows from the Lipschitz bound
for $u_i\psi_i$. If $z,w$ lie in different sets, the value of the sum
at either point is nonnegative and at most $d(z,w)$, since each bump
vanishes at the other point. Their difference is therefore at most
$d(z,w)$ in absolute value. The same argument covers points outside
all the sets.
\end{proof}

\subsection{Retaining a scale for pairs in a common bag}

There can be many distance scales in the graph, but a bag with at most
$k$ vertices uses only a bounded number. We keep one reserved pair at
each scale used by the current bag and release it when that scale is
no longer needed. A scale first appearing at a child is launched at
its parent, where its perturbation preserves the recorded parent law.

\begin{samepage}
\begin{proposition}[The co-bag embedding]\label{prop:cobag}
For every rooted tree decomposition with nonempty bags of size at most $k$, there is
an $L_1$ pseudometric $D_{\mathrm{bag}}$ on $V$ and constants
$U_{\mathrm{bag}}<\infty$ and $a_{\mathrm{bag}}>0$, depending only on
$k$, such that
\[
  D_{\mathrm{bag}}(u,v)\leq U_{\mathrm{bag}}d(u,v)
       \quad(u,v\in V),
\]
and
\[
  D_{\mathrm{bag}}(u,v)\geq a_{\mathrm{bag}}d(u,v)
       \quad\text{whenever $u,v$ belong to a common bag}.
\]
One may take
\[
  R=4\binom{k}{2},\qquad N=2R+k!+2,\qquad
  U_{\mathrm{bag}}=3N(N-1),\qquad
  a_{\mathrm{bag}}=\frac{\pi\alpha}{64}.
\]
\end{proposition}
\end{samepage}

\begin{proof}
For a bag $B$, let $\mathcal A_B$ be the set of scales $s=2^j$ for
which some distinct $u,v\in B$ satisfy
\[
  s/2\leq d(u,v)\leq4s.
\]
Call these scales active at $B$. A fixed positive distance can belong to
$[s/2,4s]$ for at most four powers of two $s$: equivalently,
$s\in[d(u,v)/4,2d(u,v)]$, an interval with endpoint ratio eight.
Consequently $|\mathcal A_B|\leq R$.

The activity interval $[s/2,4s]$ is wider than the interval $[s,2s]$
used to choose a scale for a tested pair. This slack has a geometric purpose:
if a pair at distance between $s$ and $2s$ occurs later in a consecutive run of bags where $s$ is active, at least one endpoint must lie away from the inactive launch bag.
A padded bump can then have positive amplitude at that endpoint while
vanishing at the other. Retaining its pair carries that separation to
the bag containing the tested vertices.

We use the particle construction of
Proposition~\ref{prop:particle-realization}, reserving one pair of
particles for each active scale. Initially all $N$ particles vanish.
At the root, launch a pair for each active scale using an empty interface,
and then measure the root trace. For a child of a bag $A$, make a separate
copy of the parent's list. First drop the reservations of all scales
that are inactive at this child; retain those active at both bags.
Then launch the child's newly active scales, with every update performed
at $A$, before moving to the child. This order keeps at most $R$ pairs
reserved at any time. It also leaves more than $k!$ particles free whenever
an update is required, by the choice of $N$. Different children use their
own copies. A retained pair is unchanged along every consecutive run of
bags where its scale is active.

Here is the perturbation used for a new scale $s$. Put $r=s/16$, and
sample a partition $\mathcal P$ at scale $r$ from
Theorem~\ref{thm:padded}. For $C\in\mathcal P$ give an independent
uniform coefficient $U_C\in[0,1]$ to the bump
\[
  \psi_C(z)=
  \begin{cases}
    \min\{d(z,A),d(z,V\setminus C),r\},&z\in C,\\
    0,&z\notin C.
  \end{cases}
  \qquad
  g(z)=\sum_{C\in\mathcal P}U_C\psi_C(z).
\]
For a root launch with empty interface, omit the term $d(z,A)$.
All these choices are fresh. Lemma~\ref{lem:disjoint-bumps} shows that
$g$ is $1$-Lipschitz.
It vanishes on the launch bag when that bag is nonempty. Since the scale
is active at some bag, $s\leq2\Delta$, and hence
$0\leq g\leq r\leq\Delta/8$.

At each launch use two free particles with a common weak sorting order on
the interface, replace them by $f+g$ and $f-g$ as in
Lemma~\ref{lem:trace-update}, and reserve them. An empty interface
imposes no sorting constraint. All updates have
$\operatorname{Lip}(g)\leq3$ and $|g|\leq3\Delta$.
Proposition~\ref{prop:particle-realization} therefore gives consistent
bag measures and a Markov-glued cut metric with upper bound
$3N(N-1)d$. Average over all the preparatory randomness to obtain
$D_{\mathrm{bag}}$; it is still an $L_1$ pseudometric with that upper
bound. Dropping a reservation does not delete a particle, and all updates
at a parent preserve its full trace, so this procedure also covers
branching and the ending of active runs.

It remains to prove the lower bound. Fix distinct $u,v$ in a common bag
$B$, and choose a dyadic scale $s$ with $s\leq d(u,v)<2s$.
This scale is active at $B$. Trace its reserved pair backwards to the
start of its active run. Either this was a root launch with empty
interface, or it was launched at a preceding bag $A$ where $s$ was
inactive.

In the latter case at least one of $u,v$, say $w$, has
$d(w,A)\geq s/8$. Indeed, if both distances were less than $s/8$,
nearest points $a,b\in A$ would satisfy
\[
  \tfrac34s<d(a,b)<\tfrac94s,
\]
by the triangle inequality and $s\leq d(u,v)<2s$.
In particular $a\ne b$ and $s$ would be active at $A$, a contradiction.
For an empty-interface launch take either endpoint as $w$.

Let $w'$ be the other endpoint. In the partition sampled at this launch,
$w,w'$ lie in different parts, since $d(w,w')\geq s>r$.
With probability at least $\pi$, the part $C$ containing $w$ contains
$B_d(w,\alpha r)$. On this event its bump amplitude
$a=\psi_C(w)$ is at least $\alpha r$: the distance-to-$A$ cap, when
present, is at least $s/8=2r$, and the other caps are also at least
$\alpha r$. Conditional on the partition and every coefficient other
than $U_C$, the value $g(w')$ is fixed and $g(w)=aU_C$. For every real
$b$ and $a>0$,
\[
  \int_0^1|at-b|\,dt\geq a/4.
\]
Consequently
\[
  \mathbb E\,|g(u)-g(v)|\geq\frac{\pi\alpha r}{4}
      =\frac{\pi\alpha s}{64}.
\]

The reserved particles still supply their two sweep measures in the trace
on $B$. By Equation~\eqref{eq:paired-separation}, their separation of $u,v$ is at least
$2|g(u)-g(v)|$, irrespective of the launch average $f$.
The glued metric has the prescribed trace on $B$, and all other cut
measures are nonnegative. Therefore
\[
  D_{\mathrm{bag}}(u,v)
    \geq\frac{\pi\alpha s}{32}
    \geq\frac{\pi\alpha}{64}\,d(u,v).
\]
For $u=v$ the required inequalities are immediate. This proves the
proposition.
\end{proof}
